%! AP Statistics — Unit 9, Lesson 9.4 — Group Activity (Answer Key)
\documentclass[10pt]{article}
\usepackage{apstats-article}
\usepackage{apstats-key}

\begin{document}

\pageheader{Unit 9, Lesson 9.4}{Group Activity --- Answer Key}
\namepartnerperiod

\vspace{0.1in}
\textbf{Setting Up a Significance Test for Slope}

\vspace{0.05in}
\textbf{Scenario 1:} A nutritionist randomly selects 40 adults and records their
daily fiber intake (grams, $x$) and LDL cholesterol (mg/dL, $y$). A residual plot
shows random scatter around zero with roughly equal spread. A histogram of residuals
is approximately bell-shaped. She wants to test whether fiber intake is linearly
related to LDL cholesterol.

\vspace{0.05in}
\textbf{Scenario 2:} A realtor randomly selects 50 homes sold in a city and records
square footage ($x$) and sale price (\$1000s, $y$). The residual plot shows random
scatter with no curve and roughly equal spread. $n = 50$ so normality is reasonable.
She wants to test whether square footage linearly predicts sale price.

\begin{tcolorbox}[colback=white, colframe=black!40,
    title=\textbf{Tier R --- Remediate},
    fonttitle=\bfseries, arc=2mm, left=3mm, right=3mm, top=2mm, bottom=2mm]

\textbf{Use Scenario 1 only.}

\begin{enumerate}[label=\alph*.]

  \item Name the appropriate test: \ans{$t$-test for slope}

  \item Define $\beta$: ``the true slope of the regression line relating
        \ans{fiber intake} to \ans{LDL cholesterol} for \ans{adults}.''

  \item H$_0$: $\beta = $ \ans{0} \qquad H$_a$: $\beta$ \ans{$\ne$} 0

  \item LINER conditions:

        \begin{tabularx}{\linewidth}{lXccc}
            \textbf{Letter} & \textbf{Condition} & \textbf{Met?} \\
            \hline
            L & Random scatter in residual plot & \ans{YES} \\
            I & Random sample; 10\% condition   & \ans{YES} (random sample stated; need to assume population $\ge 400$) \\
            N & Residuals bell-shaped            & \ans{YES} \\
            E & Equal spread in residual plot    & \ans{YES} \\
        \end{tabularx}

\end{enumerate}
\end{tcolorbox}

\begin{tcolorbox}[colback=white, colframe=black!40,
    title=\textbf{Tier A --- Approaching Proficiency},
    fonttitle=\bfseries, arc=2mm, left=3mm, right=3mm, top=2mm, bottom=2mm]

\begin{enumerate}[label=\alph*.]

  \item Test: \ans{$t$-test for slope} for both scenarios.

  \item \textit{Scenario 1:} $\beta = $ \ans{the true slope of the regression line
        relating LDL cholesterol to fiber intake for adults.}
        H$_0$: \ans{$\beta = 0$} \quad H$_a$: \ans{$\beta \ne 0$}

        \textit{Scenario 2:} $\beta = $ \ans{the true slope of the regression line
        relating sale price to square footage for homes in this city.}
        H$_0$: \ans{$\beta = 0$} \quad H$_a$: \ans{$\beta \ne 0$}

  \item \ans{L: Met for both --- residual plots show random scatter with no curve.
        I: Met for both --- both are random samples (assume populations $\ge 400$ and
        $\ge 500$ respectively for 10\% condition).
        N: Met for Scenario 1 (bell-shaped residuals); met for Scenario 2 ($n = 50 > 30$).
        E: Met for both --- residual plots show roughly equal spread.
        R: Met for both --- random samples stated.}

  \item \ans{The 10\% condition for I cannot be fully verified without knowing the
        total population size for each study.}

\end{enumerate}
\end{tcolorbox}

\begin{tcolorbox}[colback=white, colframe=black!40,
    title=\textbf{Tier E --- Extension},
    fonttitle=\bfseries, arc=2mm, left=3mm, right=3mm, top=2mm, bottom=2mm]

\begin{enumerate}[label=\alph*., start=5]

  \item \ans{One-sided alternative: H$_a$: $\beta > 0$ (larger homes have higher
        sale prices). This is appropriate only if she had a strong prior reason to
        believe the relationship is positive before seeing the data.}

  \item \ans{A two-sided alternative is more conservative and guards against
        data snooping --- choosing a one-sided test after observing the direction
        of $b$ inflates the Type I error rate. Unless there is a strong theoretical
        reason to rule out the other direction in advance, two-sided is preferred.}

\end{enumerate}
\end{tcolorbox}

\end{document}
